\documentclass[11pt,a4paper]{article}

%% PACHETE MATEMATICE
\usepackage{amsmath,amssymb,amsfonts}
\usepackage{amsthm}
\usepackage{mathpazo}
\usepackage{eulervm}
\usepackage{enumerate}
\usepackage{color}
\usepackage{graphicx}
\usepackage[all]{xy}
\usepackage{xcolor}
\usepackage{framed}
\usepackage[unicode]{hyperref}
\setcounter{MaxMatrixCols}{20}
\usepackage{multicol}
% \usepackage[romanian]{babel}


%% ASPECT
\linespread{1.05}
\usepackage[T1]{fontenc}
\usepackage[utf8x]{inputenc}
\usepackage[margin=2cm]{geometry}
% \usepackage[color]{showkeys}
\usepackage{anyfontsize}
\usepackage{titlesec}
\titleformat*{\section}{\large\bfseries}

\usepackage{fancyhdr}
\pagestyle{fancy}
\rhead{\small \color{gray} Notițe de Adrian Manea}
\lhead{\small \color{gray} ETTI-AM2, 2017---2018, M. Joița \& A. Niță}


\usepackage[autostyle = false, style = german]{csquotes}
\usepackage[romanian]{babel}
\newcommand\qq{\enquote}

%% COMENZILE MELE
\renewcommand*{\proofname}{Dem.:}
\newcommand\bb{\mathbb}
\newcommand\dr{\mathrm}
\newcommand\kal{\mathcal}
\newcommand\fk{\mathfrak}
\newcommand\op{\oplus}
\newcommand\ot{\otimes}
\newcommand\ds{\displaystyle}
\newcommand\id{\indent}
\newcommand\nid{\noindent}
\newcommand\seq{\subseteq}
\newcommand\sneq{\subsetneq}
\newcommand\speq{\supseteq}
\newcommand\spneq{\supsetneq}
\newcommand\rar{\rightarrow}
\newcommand\Rar{\Rightarrow}
\newcommand\xrar{\xrightarrow}
\newcommand\lar{\leftarrow}
\newcommand\Lar{\Leftarrow}
\newcommand\xlar{\xleftarrow}
\newcommand\lrar{\leftrightarrow}
\newcommand\Lrar{\Leftrightarrow}
\newcommand\ideal{\trianglelefteq}
\newcommand\ai{\text{ a.\^{i}. }}
\newcommand\sk{(\textit{Schi\c{t}\u{a}}) }
\newcommand\rhup{\rightharpoonup}
\newcommand\rhdn{\rightharpoondown}
\newcommand\lhup{\leftharpoonup}
\newcommand\lhdn{\leftharpoondown}
\newcommand\vid{\emptyset}
\newcommand\wh{\widehat}
\newcommand\ol{\overline}
\newcommand\NN{\mathbb{N}}
\newcommand\RR{\mathbb{R}}
\newcommand\ZZ{\mathbb{Z}}
\newcommand\QQ{\mathbb{Q}}
\newcommand\CC{\mathbb{C}}

%% STILURILE MELE
\newtheoremstyle{dotless}{}{}{\itshape}{}{\bfseries}{:}{ }{}

\theoremstyle{dotless}
\newtheorem{theorem}{Teorem\u{a}}[section]
\newtheorem{proposition}{Propozi\c{t}ie}[section]
\newtheorem{lemma}{Lem\u{a}}[section]
\newtheorem{corollary}{Corolar}[section]
\newtheorem{exercise}{Exerci\c{t}iu}

\newtheoremstyle{dotlessdr}{}{}{}{}{\bfseries}{:}{ }{}
\theoremstyle{dotlessdr}
\newtheorem{example}{Exemplu}[section]
\newtheorem{definition}{Defini\c{t}ie}[section]
\newtheorem{remark}{Observa\c{t}ie}[section]




\begin{document}
% \definecolor{labelkey}{RGB}{222,0,0}

\begin{center}
  {\large\textbf{Seminar 4 \\
      Ecuații și sisteme autonome. \\
      Stabilitate. Integrale prime}}
\end{center}

\vspace{2cm}

\section{Ecuații autonome (nu conțin pe $x$)}

În cazul acestor ecuații, putem micșora ordinul cu o unitate dacă notăm $y' = p$
și luăm pe $y$ variabilă independentă.

\begin{remark}\label{rk:autonom-pierd}
  Există posibilitatea să pierdem soluții de forma $y = c$ prin această metodă,
  deci trebuie verificat ulterior dacă a fost cazul.
\end{remark}

Exemplu 1: $1 + {y'}^2 = 2yy^{\prime\prime}$.

\emph{Soluție:} Luăm $y' = p$ drept funcție și pe $y$ drept variabilă independentă.
Obținem:
\[
  y^{\prime\prime} = \frac{d}{dx} \big( \frac{dy}{dx} \big) = p \frac{dp}{dy}.
\]

Atunci ecuația devine $\frac{2pdp}{1 + p^2} = \frac{dy}{y}$, ce are drept soluție
$y = c_1(1 + p^2)$.

Acum trebuie să obținem pe $x$ ca funcție de $p$ și $c_1$. Deoarece $dx = \frac{1}{p}dy$,
iar $dy = 2c_1 pdp$, rezultă $dx = 2c_1 dp$. Așadar, $x = 2c_1 p + c_2$, iar soluția
generală este $x(p) = 2c_1 p + c_2$. Cum $y(p) = c_1(1 + p^2)$, rezultă:
\[
  y = c_1 + \frac{(x - c_2)^2}{4c_1}.
\]

\vspace{1cm}

Exemplu 2: $y^{\prime\prime} + {y'}^2 = 2e^{-y}$.

\emph{Soluție:} Notăm $y' = p$ și luăm ca necunoscută $p = p(y)$,
de unde obținem:
\[
  y^{\prime\prime} = \frac{dy'}{dx} = p'p \Rightarrow p'p + p^2 = 2e^{-y}.
\]
Dacă notăm $p^2 = z$, obținem o ecuație liniară neomogenă:
\[
  z' + 2z = 4e^{-y},
\]
ce are ca soluție generală $z(y) = c_1 e^{-2y} + 4e^{-y}$. Revenind la $y$,
avem:
\[
  z = p^2 = {y'}^2 \Rightarrow y' = \pm \sqrt{c_1 e^{-2y} + 4e^{-y}},
\]
adică ecuația:
\[
  \frac{dy}{\pm\sqrt{c_1 e^{2y} + 4e^{-y}}} = dx,
\]
ce are ca soluție: $x + c_2 = \pm \frac{1}{2} \sqrt{c_1 + 4e^y}$. Echivalent, în forma
implicită:
\[
  e^y + \frac{c_1}{4} = (x + c_2)^2.
\]

\vspace{1cm}

Similar putem proceda și în cazul:
\[
  y^{\prime\prime\prime} + y^{\prime\prime} = \sin x.
\]

\emph{Indicație:} Putem nota $y^{\prime\prime} = z$ și atunci ecuația devine:
\[
  z' + z = \sin x,
\]
care este o ecuație liniară neomogenă, cu soluția:
\[
  z(x) = e^{-x}(c_1 + \frac{e^x}{2}(\sin x - \cos x)).
\]
Mai departe, integrăm succesiv:
\begin{align*}
  y'(x) &= \int z(x) dx + c_2 = -c_1 e^{-x} - \frac{1}{2} \cos x - \frac{1}{2}\sin x + c_2 \\
  y(x) &= \int y'(x) dx + c_3 = c_1 e^{-x} - \frac{1}{2} \sin x + \frac{1}{2} \cos x + c_2 x + c_3.
\end{align*}

%%%%%%%%%%%%%%%%%%%%%%%%%%%%%%%%%%%%%%%%%%%%%%%%%%%%%%%%%%%%%%%%%%%%%%


\section{Sisteme diferențiale autonome. Integrale prime}
\begin{definition}\label{def:s-d-autonom}
  Un sistem diferențial de forma:
  \[
    x'_j = v_j(x_1, \dots, x_n), \quad j = 1, 2, \dots, n,
  \]
  unde $v = (v_1, \dots, v_n) : U \to \RR^n$ este un cîmp de vectori de
  clasă $\kal{C}^r, r \geq 1$, definit într-un domeniu $U \seq \RR^n$
  se numește \emph{sistem diferențial autonom}.
\end{definition}

Sistemul de mai sus poate fi scris și într-o \emph{formă simetrică}:
\[
  \frac{dx_1}{v_1} = \dots = \frac{dx_n}{v_n}.
\]

Dacă $f : U \to \RR$ este o funcție de clasă $\kal{C}^1(U)$, atunci pentru orice
$x \in U$ se poate considera \emph{derivata lui $f$ în $x$ și în direcția vectorului %
  $v(x)$}, notată $\frac{df}{dv}(x)$, definită prin formula cunoscută:
\[
  \frac{df}{dv}(x) = \sum_{i = 1}^n \frac{\partial f}{\partial x_i} v_i(x).
\]

\begin{definition}\label{def:int-prim}
  Fie $v : U \to \RR^n$ un cîmp de vectori și $f : U \to \RR$ o funcție de clasă
  $\kal{C}^1(U)$. Funcția se numește o \emph{integrală primă} a sistemului diferențial
  autonom $x' = v(x), x \in U$ dacă derivata sa în direcția cîmpului de vectori $v$
  este nulă în fiecare punct din $U$, adică $\frac{df}{dv} = 0$.
\end{definition}

Echivalent, putem formula definiția astfel: $f : U \to \RR$ este o integrală
primă pentru sistemul diferențial autonom $x' = v(x)$ dacă și numai dacă oricare
ar fi soluția $x = \varphi(t), \varphi : I \to U$, funcția $f \circ \varphi$ este
constantă pe $I$. Uneori, mai putem scrie pe scurt $f(x_1, \dots, x_n) = c$, constant.
De aceea, putem spune că integralele prime reprezintă \emph{legi de conservare}.

Tehnic, pentru a găsi integralele prime asociate unui sistem autonom, se scrie
sistemul în forma simetrică, iar apoi, folosind proprietățile rapoartelor egale,
se ajunge la un raport de forma $\frac{df}{0}$, egal cu rapoartele precedente.
Atunci $f$ va fi integrală primă, deoarece va rezulta $df = 0$, adică $f$ constantă
în lungul curbelor integrale.

\textbf{Exemplu:} Să se găsească integralele prime ale sistemului simetric:
\[
  \frac{dx}{z - y} = \frac{dy}{x - z} = \frac{dz}{y - x}.
\]

\emph{Soluție:} Folosind proprietățile proporțiilor, putem scrie sistemul:
\[
  \frac{dx}{z - y} = \frac{dy}{x - z} = \frac{dz}{y - x} = %
  \frac{dx + dy + dz}{z - x + x - z + y - x} = \frac{dx + dy + dz}{0}
\]

De asemenea, mai putem obține o integrală primă astfel:
\[
  \frac{dx}{z - y} = \frac{dy}{x - z} = \frac{dz}{y - x} = %
  \frac{xdx + ydy + zdz}{x(z - x) + y(x - z) + z(y - x)} = %
  \frac{xdx + ydy + zdz}{0}.
\]

Așadar, avem:
\begin{align*}
  dx + dy + dz &= 0 \\
  xdx + ydy + zdz &= 0,
\end{align*}
de unde obținem două integrale prime:
\[
  x + y + z = c_1, \quad x^2 + y^2 + z^2 = c_2.
\]

\begin{remark}\label{rk:int-prima-t}
  Dacă avem un sistem diferențial neautonom, adică $x' = v(x, t)$, cu $t \in \RR$
  și $x \in U$, atunci o integrală primă a sistemului va fi o funcție
  $f : U \times \RR \to \RR$, diferențiabilă și \emph{dependentă de timp}. Ea
  se obține din sistemul anterior, adăugînd ecuația $t' = 1$.
\end{remark}

%%%%%%%%%%%%%%%%%%%%%%%%%%%%%%%%%%%%%%%%%%%%%%%%%%%%%%%%%%%%%%%%%%%%%%

\section{Stabilitatea soluțiilor sistemelor diferențiale}

Considerăm un sistem diferențial $x' = v(t, x)$. Presupunem că sînt îndeplinite
condițiile teoremei fundamentale de existență și unicitate a soluției problemei
Cauchy pentru $t \in [t_0, \infty)$ și $x \in U, U \seq \RR^n$ este un deschis.
Așadar, pentru orice $x_0 \in U$, există și este unică o soluție $x = \varphi(t)$,
cu $\varphi : [t_0, \infty) \to U$, astfel încît $\varphi(t_0) = x_0$.

\begin{definition}\label{def:stabila-PL}
  O soluție $x = \varphi(t), \varphi : [t_0, \infty) \to U$ se numește \emph{stabilă %
    spre $\infty$ în sens Poincar\'{e}-Liapunov} sau, echivalent, $x_0 = \varphi(t_0)$
  se numește \emph{poziție de echilibru} dacă la variații mici ale lui $x_0$ obținem
  variații mici ale soluției. Formal, pentru orice $\varepsilon > 0$, există
  $\delta(\varepsilon) > 0$ astfel încît pentru $\widetilde{x_0}\in\RR^n$ cu
  $||\widetilde{x}_0 - x_0|| < \delta(\varepsilon)$, soluțiile $\widetilde{\varphi}(t)$
  și $\varphi(t)$ satisfac inegalitatea $|\widetilde{\varphi}(t) - \varphi(t)| < \varepsilon$,
  pentru orice $t \in [t_0, \infty)$.
\end{definition}

\begin{definition}\label{def:stabila-0}
  Poziția de echilibru $x = 0$ se numește \emph{stabilă în sens Poincar\'{e}-Liapunov}
  dacă pentru orice $\varepsilon > 0$, există $\delta > 0$, care depinde doar de
  $\varepsilon$ astfel încît pentru orice $x_0 \in U$ pentru care $||x_0|| < \delta$,
  soluția $\varphi$ a sistemului cu condiția inițială $\varphi(0) = x_0$ se prelungește
  pe întreaga semiaxă $t > 0$ și satisface inegalitatea $||\varphi(t)|| < \varepsilon$
  pentru orice $t > 0$.
\end{definition}

\begin{definition}\label{def:as-stabila}
  Poziția de echilibru $x = 0$ a sistemului diferențial autonom se numește
  \emph{asimptotic stabilă} dacă este stabilă și, în plus, pentru soluția $\varphi(t)$
  din definiția de mai sus, avem:
  \[
    \lim_{t \to \infty} \varphi(t) = 0.
  \]
\end{definition}

În cazul sistemelor diferențiale liniare și omogene, de forma:
\[
  x' = Ax, \quad A : \RR^n \to \RR^n, x \in \RR^n
\]
presupunem că $A$ este izomorfism (echivalent, matricea $A$ este inversabilă).
Atunci $x = 0$ este singurul punct singular al cîmpului $v(x) = Ax$, deci $x = 0$
este singura poziție de echilibru a sistemului.

În exerciții, stabilitatea se va studia folosind următoarea:

\begin{theorem}[Poincar\'{e}-Liapunov]\label{thm:stab-pl}
  Păstrînd contextul și notațiile de mai sus, dacă toate valorile proprii ale
  matricei $A$ au partea reală negativă, atunci poziția de echilibru $x = 0$ este
  asimptotic stabilă.

  Dacă există $\lambda \in \sigma(A)$ cu $\mathrm{Re}\lambda > 0$, atunci $x = 0$
  este instabilă.
\end{theorem}

%%%%%%%%%%%%%%%%%%%%%%%%%%%%%%%%%%%%%%%%%%%%%%%%%%%%%%%%%%%%%%%%%%%%%%

\section{Exerciții}

1. Studiați stabilitatea poziției de echilibru $x = 0$ pentru sistemele diferențiale:
\begin{enumerate}[(a)]
\item $\begin{cases}
    x' &= -4x + y \\
    y' &= -x - 2y
  \end{cases}$;
\item $\begin{cases}
    x' &= -x + z \\
    y' &= -2y - z \\
    z' &= y - z
  \end{cases}$;
\item $\begin{cases}
    x' &= 2y \\
    y' &= x + a
  \end{cases}, a \in \RR$.
\end{enumerate}

\vspace{1cm}

2. Să se afle pentru ce valori ale lui $a \in \RR$ soluția nulă a sistemului
de mai jos este asimptotic stabilă:
\[
  \begin{cases}
    x' &= ax + y + z^2 \\
    y' &= (2 + a)x + ay + \cos y - 1 \\
    z' &= x + y - z
  \end{cases},
\]
unde derivatele sînt considerate în raport cu variabila $t$.

\emph{Indicație:} Considerăm sistemul liniarizat:
\[
  \begin{cases}
    x' &= ax + y \\
    y' &= (2 + a)x + ay \\
    z' &= x + y - z
  \end{cases}
\]

\vspace{1cm}

3. Să se determine liniile de cîmp pentru cîmpurile vectoriale de pe $\RR^3$:
\begin{enumerate}[(a)]
\item $\vec{v} = x\vec{i} + y \vec{j} + xyz \vec{k}$;
\item $\vec{v} = (2z - 3y)\vec{i} + (6x - 2z)\vec{j} + (3y - 6x)\vec{k}$;
\item $\vec{v} = (xz + y) \vec{i} + (x + yz) \vec{j} + (1 - z^2)\vec{k}$;
\item $\vec{v} = (x^2 + y^2)\vec{i} + 2xy \vec{j} - z^2 \vec{k}$.
\end{enumerate}

\emph{Indicație (a):} Scriem sistemul autonom asociat cîmpului vectorial $\vec{v}$
sub forma:
\[
  \frac{dx}{x} = \frac{dy}{y} = \frac{dz}{xyz}.
\]
Din prima egalitate, rezultă $x = c_1 y$, deci a doua egalitate devine:
\[
  c_1 ydy = \frac{dz}{z} \Rightarrow c_1 \frac{y^2}{2} = \ln |z| + c_2.
\]
Așadar, obținem $\frac{x}{y} \cdot \frac{y^2}{2} = \ln |z| + c_2$.

Rezultă că liniile de cîmp pentru $\vec{v}$ sînt curbele:
\[
  \begin{cases}
    \frac{x}{y} &= c_1 \\
    xy - \ln|z| &= c_2.
  \end{cases}
\]

\vspace{1cm}

(b) Sistemul poate fi scris sub forma:
\[
  \frac{dx}{2z - 3y} = \frac{dy}{6x - 2z} = \frac{dz}{3y-6x} = \frac{dx + dy + dz}{0}.
\]
Așadar, $dx + dy + dz = 0$, deci $x + y + z = c_1$.

Din forma inițială obținem și:
\[
  \frac{3xdx}{3x(2z - 3y)} = \frac{\frac{3}{2}ydy}{\frac{3}{2}y(3x - z)} = \frac{zdz}{z(y - 2x)} = %
  \frac{3xdx + \frac{3}{2}ydy + zdz}{0},
\]
deci $3xdx + \frac{3}{2}ydy + zdz = 0$, adică
$3\frac{x^2}{2} + \frac{3}{4}y^2 + \frac{z^2}{22} = c_2$.

\vspace{1cm}

(c) Obținem:
\[
  \frac{dx + dy}{(x + y)(z + 1)} = \frac{dz}{1 - z^2} \Rightarrow %
  \frac{d(x + y)}{x + y} = -\frac{dz}{z - 1}.
\]
Rezultă $\ln|x + y| + \ln|z - 1| = \ln c_1$, adică $(x + y)(z - 1) = c_1$.

Apoi:
\[
  \frac{dy - dx}{x + zy - xz - y} = \frac{dy - dx}{(x - y)(1 - z)} = \frac{dz}{1 - z^2}.
\]

Rezultă $-\frac{d(x - y)}{x - y} = \frac{dz}{z + 1}$, de unde $(x - y)(z + 1) = c_2$.

Așadar, liniile de cîmp sînt curbele:
\[
  \begin{cases}
    (x + y)(z - 1) &= c_1 \\
    (x - y)(z + 1) &= c_2
  \end{cases}
\]

\vspace{1cm}

(d) Sistemul simetric rezultat este:
\[
  \frac{dx}{x^2 + y^2} = \frac{dy}{2xy} = \frac{dz}{-z^2}.
\]

Prin adunarea primelor două rapoarte, obținem:
\[
  \frac{d(x + y)}{(x + y)^2} = \frac{dz}{-z^2},
\]
adică $\frac{1}{x + y} + \frac{1}{z} = c_1$.

Prin scădere, avem $\frac{d(x - y)}{(x - y)^2} = \frac{dz}{-z^2}$, deci
$\frac{1}{x - y} + \frac{1}{z} = c_2$.

Liniile de cîmp se obțin:
\[
  \begin{cases}
    \frac{1}{x + y} + \frac{1}{z} &= c_1 \\
    \frac{1}{x - y} + \frac{1}{z} &= c_2
  \end{cases}
\]

%%%%%%%%%%%%%%%%%%%%%%%%%%%%%%%%%%%%%%%%%%%%%%%%%%%%%%%%%%%%%%%%%%%%%%


\end{document}